% tex/40_termination.tex
\section{Theorem: termination via well-founded rank}
\label{sec:termination}

This section records the ``general theorem C'' as a finished Lean artifact (C2) and shows visible reuse (C3).

\subsection{Well-foundedness from a rank into \texorpdfstring{$\mathbb{N}$}{N}}
Let \((X,\rank)\) be a ranked object with \(\rank : X\to\mathbb{N}\). Define the descending relation
\[
  \mathrm{Desc}(x,y) :\!\!\iff \rank(x) < \rank(y).
\]

\begin{theorem}[C2: \(\mathrm{Desc}\) is well-founded]
\label{thm:wf-desc}
\(\mathrm{Desc}\) is well-founded.
\end{theorem}

\paragraph{Proof sketch (one line).}
Use the \wfmeasure\ induced by \(\rank : X\to\mathbb{N}\):
\[
  \text{\leaninline{by simpa [Desc] using (measure rank).wf}.}
\]
\emph{Terminology warning:} this is not measure theory (\code{MeasureTheory}); it is the standard
well-founded construction from a function into \(\mathbb{N}\).

\subsection{Two applications (C3)}
Keep these short and visibly reusing the kernel.

\begin{example}[List peeling terminates]
Define \(\mathrm{TailRel}(x,y) :\!\!\iff \exists a,\ y=a::x\). Then \(\mathrm{TailRel}\) is well-founded.
\end{example}

\begin{example}[Absolute-value descent on \(\mathbb{Z}\)]
Define \(\mathrm{AbsDesc}(x,y) :\!\!\iff |x|<|y|\) with \( |x| = \mathrm{Int.natAbs}(x)\).
Then \(\mathrm{AbsDesc}\) is well-founded.
\end{example}

\subsection{Where to paste T4--T5}
Paste (or paraphrase) the frozen T4/T5 text here:
\begin{itemize}
\item the C1/C2 statement and the measure proof sketch,
\item the two C3 examples (as above),
\item the contribution statement (allowed from T5).
\end{itemize}
